\corrsection{10}\label{C10-CORR}
\textbf{Usage professeur.} Les réponses qui suivent concernent les tâches non résolues dans le cours. Faire d’abord expliciter la démarche ; fournir une aide graduée avant la solution. Les durées sont des estimations à étalonner en classe.
\subsection*{Diagnostic et remédiation}
\begin{enumerate}
\item $2\pi$. \textit{Action :} Relier rayon, diamètre et périmètre.
\item $90^\circ$. \textit{Action :} Partager un cercle en quatre secteurs.
\item $\dfrac{\sqrt 2}{2}$. \textit{Action :} Rationaliser et vérifier le carré.
\item $(a;-b)$. \textit{Action :} Tracer les deux points.
\item $-\dfrac{\pi}{3}$. \textit{Action :} Mettre au même dénominateur.
\item $\dfrac{3}{5}$. \textit{Action :} Identifier l’angle de référence.
\end{enumerate}
\subsection*{Activité de recherche}
Triangle équilatéral puis triangle rectangle isocèle. Pythagore donne respectivement hauteur $\dfrac{\sqrt 3}{2}$ et coordonnées égales $\dfrac{\sqrt 2}{2}$. Les signes sont positifs dans le premier quadrant.
\subsection*{Applications immédiates}
\begin{enumerate}
\item $\dfrac{\pi}{4}$, $\dfrac{2\pi}{3}$ ; longueur $\dfrac{3\pi}{2}$ m.
\item $\dfrac{7\pi}{2}$ a l’image de $\dfrac{3\pi}{2}$ ; $-\dfrac{5\pi}{3}$ celle de $\dfrac{\pi}{3}$.
\item $\sin ^2x=\dfrac{16}{25}$ ; le sinus est négatif, donc $-\dfrac{4}{5}$.
\item $(-\dfrac{1}{2};-\dfrac{\sqrt 3}{2})$ et $(\dfrac{\sqrt 2}{2};-\dfrac{\sqrt 2}{2})$.
\item $h=\dfrac{\sqrt 3}{2}$ ; nouvelle corde $\dfrac{1}{\sqrt {2+\sqrt3}}=\sqrt {2-\sqrt3}$.
\end{enumerate}
\subsection*{Exercices : démarche et résultats}
\paragraph{Exercice 10.1 -- Convertir et mesurer.} $\dfrac{5\pi}{6}$ et $-\dfrac{\pi}{2}$ ; arc $\dfrac{4\pi}{3}$ cm. Le signe négatif indique une orientation ; une longueur géométrique est positive.
\paragraph{Exercice 10.2 -- Réduire les angles.} $\dfrac{5\pi}{6}$, $\dfrac{\pi}{4}$, $\pi$. Retirer ou ajouter des multiples de $2\pi$.
\paragraph{Exercice 10.3 -- Valeurs exactes.} $(-\dfrac{1}{2};\dfrac{\sqrt 3}{2})$ ; $(\dfrac{\sqrt 3}{2};-\dfrac{1}{2})$ ; $(-\dfrac{\sqrt 2}{2};-\dfrac{\sqrt 2}{2})$, sous forme (cosinus ; sinus).
\paragraph{Exercice 10.4 -- Retrouver le signe.} $\cos ^2x=1-\dfrac{25}{169}=\dfrac{144}{169}$ ; cosinus négatif, donc $-\dfrac{12}{13}$.
\paragraph{Exercice 10.5 -- Une roue.} Avancement $r\theta =0{,}3(\dfrac{5\pi}{3})=\dfrac{\pi}{2}$ m. Cela correspond au déroulement de l’arc ; la corde relie deux positions sur le cercle et n’a pas cette longueur. Le point matériel sur la roue décrit une trajectoire différente par rapport au sol.
\paragraph{Exercice 10.6 -- Corde et produit scalaire.} $A(1;0)$, $B(-\dfrac{1}{2};\dfrac{\sqrt 3}{2})$ ; $AB^2=\dfrac{9}{4}+\dfrac{3}{4}=3$. Al-Kashi : $1+1-2\cos (\dfrac{2\pi}{3})=3$. Donc $AB=\sqrt3$.
\paragraph{Exercice 10.7 -- Lire une équation sur le cercle.} La verticale d’abscisse $\dfrac{1}{2}$ coupe le cercle aux angles $\dfrac{\pi}{3}$ et $\dfrac{5\pi}{3}$. Les deux points ont des sinus opposés.
\paragraph{Exercice 10.8 -- Reprise différée.} $\dfrac{25\pi}{4}=6\pi +\dfrac{\pi}{4}$. Trois tours complets ne changent pas l’image, donc les deux valeurs valent $\dfrac{\sqrt 2}{2}$.
\paragraph{CQFD, exercice 32 p. 220, questions a et b.} a : $\cos(\dfrac{6\pi}5)=\cos(\pi+\dfrac\pi5)=-\dfrac{1+\sqrt5}{4}$. b : $\cos(-\dfrac\pi5)=\cos(\dfrac\pi5)=\dfrac{1+\sqrt5}{4}$, par parité du cosinus.
\paragraph{CQFD, exercice 39 p. 220, questions 1, 2, 3 et 5.} 1 : $\dfrac\pi6+\dfrac\pi3=\dfrac\pi2$. 2 : oui, $\cos(\dfrac\pi3)=\dfrac12=0{,}5$ est décimal. 3 : la mesure principale dans $]-\pi;\pi]$ est $-\dfrac{5\pi}6$. 5 : $\sin(\dfrac{4\pi}3)=-\dfrac{\sqrt3}2$.
\subsection*{QCM : réponses et raisons}
\textbf{QCM 1.}
\begin{enumerate}
\item B. Demi-tour.
\item A. Triangle équilatéral.
\item D. Point inférieur du cercle.
\item C. Équation du cercle unité.
\item B. Tours complets.
\end{enumerate}
\textbf{QCM 2.}
\begin{enumerate}
\item C. Symétrie verticale.
\item A. Triangle rectangle isocèle.
\item D. Trois quarts de tour.
\item B. $r\theta=6$.
\item C. Abscisse négative à l’intérieur du quadrant.
\end{enumerate}
\subsection*{Vrai ou faux}
\begin{enumerate}
\item Faux : en 0, les deux membres seraient 1 et $-1$. Le cosinus est pair.
\item Vrai : un tour complet conserve l’image.
\item Faux : $\dfrac{\pi}{6}$ et $\dfrac{5\pi}{6}$ ont même sinus mais cosinus opposés.
\end{enumerate}
\paragraph{Différenciation.} Pour un blocage technique, proposer la première étape seulement (domaine, facteur ou univers de référence), puis retirer l’aide. Pour un élève autonome, demander une seconde méthode et le traitement d’un cas limite. À l’oral, autoriser une reformulation initiale dans la langue maîtrisée, puis exiger une phrase mathématique précise en français.
